% ============================================================
% band0/chapters/frontmatter.tex
% Didaktisches Konzept und Zielgruppe -- nur für Band 0
% ============================================================

\section*{Zwei Niveaus}

Dieses Buch arbeitet mit zwei Ebenen.

\subsection*{Der Hauptweg}
Der Hauptweg zeigt, wie ein konkretes Problem gelöst wird. Er ist praktisch,
direkt und auf Umsetzung ausgelegt. Wer schnell zu einem funktionierenden
Ergebnis kommen will, folgt dem Hauptweg.

\subsection*{Unter der Haube}
Die blau markierten Vertiefungsbereiche erklären die technischen Hintergründe:
Betriebssystemkonzepte, Architektur, Sicherheit und mathematische Denkmodelle.

Der Hauptweg funktioniert auch ohne die Vertiefung. Wer aber genauer verstehen
will, \emph{warum} eine Lösung funktioniert, findet dort die fachliche Tiefe.

\section*{Für wen ist dieses Buch gedacht?}

\begin{itemize}
  \item Einsteigerinnen und Einsteiger ohne Linux-Erfahrung,
  \item Schülerinnen und Schüler technischer Bildungsgänge,
  \item Studierende technischer Studiengänge,
  \item Maker und Praktiker,
  \item Lehrkräfte, die Robotik systematisch unterrichten möchten,
  \item alle, die später mit ROS2, Python und C++ Roboter entwickeln wollen.
\end{itemize}

\section*{Was dieses Buch nicht ist}

Dieses Buch ist kein vollständiges Linux-Handbuch und kein Docker-Referenzwerk.
Es ist eine gezielte Vorbereitung auf moderne Robotikentwicklung.
